\documentclass[10pt]{article}

\usepackage[utf8]{inputenc}

\usepackage[T1]{fontenc}

\usepackage{amsmath}

\usepackage{amsfonts}

\usepackage{amssymb}

\usepackage[version=4]{mhchem}

\usepackage{stmaryrd}

\usepackage{mathrsfs}



\begin{document}

HECTEKTPAЛЬНАЯ ОСОБEHНОСТЬ с а о сопря женного операто ра $A$, действующего в гильбертовом пространстве $H,-$ сингулярность аналитического продолжения его резольвенты



\[

P_{1}(A-\lambda I)^{-1} P_{2}

\]



окаймленной ортопроекторами $P_{i}, i=1,2$, на некоторые подпространства $H_{i}$ из $H$. В то время как естественной особенностью аналитичности самой резольвенты является спектральный лист переменной $\lambda$ - дополнение спектра $\sigma(A)$ на комплексной плоскости $\lambda,-$ окаймленная резольвента в нек-рых случаях допускает аналитич. продолжение на «нефизический лист» через разрезы, отвечающие абсолютно непрерывному спектру. Интегрируя билинейную форму резольвенты $\left\langle R_{\lambda} e_{i}, e_{k}\right\rangle, e_{i} \in H_{i}$, с экспонентой по контуру $\gamma$ вокруг спектра оператора $A$ :



\[

\left\langle e^{i A t}, e_{i}, e_{k}\right\rangle=-\frac{1}{2 \pi i} \oint_{\gamma} e^{i \lambda t}\left\langle R_{\lambda} e_{i}, e_{k}\right\rangle d \lambda,

\]



можно, пользуясь продолжимостью окаймленной резольвенты на нефизич. лист при $t>0$, продеформировать контур интегрирования, убирая его из полуплоскости нижней на верхнюю полуплоскость нефизич. листа и обходя встречающиеся на пути неспектральные особенности. Интеграл на нефизич. листе вокруг полюса $\lambda_{0}$ аналитич. продолжения дает экспоненциально убывающий вклад



\[

\text { Rese }\left.e^{i \lambda t}\left\langle R_{\lambda}, e_{i}, e_{k}\right\rangle\right|_{\lambda=\lambda_{0}},

\]



а участки деформированного контура интегрирования, примыкающие к вещественной оси $\lambda,-$ вклады, убывающие степенным образом, как $O\left(t^{-l}\right), l>0$.



В тех спектральных задачах, где существует матрица рассеяния $S(k)$, часто оказывается, что она, будучи окаймлена проекторами на нек-рые подпространства $E_{i}, i=1,2$, вспомогательного пространства также допускает аналитич. продолжение с вещественной оси спектральной переменной $k$ в комплексную плоскость. Если исходная $S$-матрица унитарна, то область аналитичности продолженной $S$-матрицы всегда можно считать симметричной в силу принципа симметрии. Корни продолженной $S$-матрицы называются р е 3 она нс а м и.



Наиболее интересной и аналитически богатой является ситуация "резонансного рассеяния", когда «неспектральные особенности» совпадают с резонансами. Этот случай подробно проанализирован П. Лаксом и Р. Филлипсом (см. Лакса-Филлипса подход). Он имеет место тогда, когда рассматриваемый самосопряженный оператор $A$ является генератором унитарной группы $U_{t}=\exp (i A t)$, обладающей ортогональными приходящими и уходящими подпространствами $D_{\mp}$. При этом резольвента оператора $A$, окаймленная проекторами на трансляционно инвариантное подпространство $K=H \ominus\left(D_{-} \oplus D_{+}\right)$, совпадает в нижней полуплоскости $\operatorname{Im} K<0$ с резольвентой генератора $B$ полугруппы сжатий $Z_{t}=\left.P_{K} e^{i A t}\right|_{K}$ :



\[

\left.P_{K}(A-\lambda I)^{-1}\right|_{K}=(B-k I)^{-1}, \operatorname{Im} k<0,

\]



и, значит, продолжается в верхнюю полуплоскость на дополнение спектра оператора $B$. При этом матрица рассеяния $S(k)$ оператора $A$ оказывается характеристич. функцией генератора $B$ и множество особенностей $S^{-1}(k)$ совпадает со спектром оператора $B$.



В этом случае главные части окаймленной резольвенты в еe полюсах (резонансах) записываются через собственные функции генератора $B$ - резонансные состояния. Анализ резонансных состояний эквивалентен спектральному анализу общего диссипативного оператора. Б. С. Павлов.



НЕСТАЦИОНАРНАЯ СВЯЗЬ - связь механическая, изменяющаяся со временем.



384 HETEP HECTPAННЫЙ TOK - см. Ток.



HETEP TEOPEMA: для всякой физической системы, уравнения движения которой могут быть получены из вариационного принципа, каждому однопараметрическому непрерывному преобразованию, оставляющему вариационный функционал инвариантным, отвечает один дифференциальный сохранения закон и, главное, позволяет явно выписать сохраняющуюся величину. Установлена в работах ученых геттингенской школы Д. Гильберта (D. Hilbert), Ф. Клейна (F. Klein) и Э. Нетер (Е. Noether). Н. т.- самое универсальное средство, позволяющее находить законы сохранения в лагранжевой классич. механике, теории поля, квантовой теории и т. д.



В классич. механике для системы с действием



\[

S=\int_{T} L\left[q_{a}(t), \dot{q}_{a}(t)\right] d t

\]



( $L$ - Лагранжа функция, зависящая от обобщенных координат $q_{a}$ и скоростей $\dot{q}_{a}$ ) инвариантность $S$ относительно образующих группу преобразований с параметром $\varepsilon$ :



\[

t \rightarrow t^{\prime}=t+\Lambda(q, t) \varepsilon, q_{a}(t) \rightarrow q_{a}^{\prime}\left(t^{\prime}\right)=q_{a}(t)+\Lambda_{a}(q, t) \varepsilon

\]



[где задающие преобразование функции $\Lambda(q, t), \Lambda_{a}(q, t)$ зависят от совокупности координат $\left\{q_{a}\right\} \equiv q$ и времени] влечет за собой, согласно Н.т., сохранение во времени величины



\[

Q=\left[L-\sum_{a} \frac{\partial L}{\partial \dot{q}_{a}} \dot{q}_{a}\right] \Lambda+\sum_{a} \frac{\partial L}{\partial q_{a}} \Lambda_{a}

\]



В частности, из инвариантности $S$ относительно (1) с $\Lambda_{a}=0$, $\Lambda=1$, то есть из однородности времени, следует закон сохранения энергии:



\[

-\mathscr{E}=L-\sum_{a} \frac{\partial L}{\partial \dot{q}_{a}} \dot{q}_{a}=\text { const }

\]



В этом случае $L$ не зависит от времени явно. Подобным же образом из инвариантности $S$ по отношению к пространственным сдвигам $\left(\Lambda=0, \Lambda_{a}=1\right)$ следует закон сохранения импульса, а из изотропии пространства - закон сохранения трехмерного момента.



В гамильтоновом описании, то есть когда $Q$ выражены через канонич. переменные - обобщенные координаты и импульсы (для простоты считаем, что явные зависимости от времени отсутствуют): 1) Пуассона скобка $Q$ с гамильтонианом $H$ равна нулю, 2) изменение любой динамич. переменной $F$ при преобразовании (1) определяется ее скобкой Пуассона с $Q$. В этом контексте утверждение Н. т. становится как бы тривиальным, следующим из одной лишь антисимметрии скобок Пуассона:



\[

0=d H / d \varepsilon=(Q, H) \rightarrow 0=(H, Q)=d Q / d t .

\]



Если преобразования симметрии образуют не однопараметрич. группу, то между $Q_{A}$ должны выполняться соотношения в скобках Пуассона, воспроизводящие алгебру Ли генераторов соответствующей группы. Так, напр., три компоненты момента должны удовлетворять соотношению в скобках Пуассона



\[

\left(M_{i} M_{k}\right)=-e_{i} k_{l} M_{l}, i, k, l=1,2,3

\]



(где $e_{i} k_{l}$ - символ Леви-Чивиты), воспроизводящему алгебру Ли группы трехмерных вращений $O(3)$.



Особо важное значение Н. т. приобретает в квантовой теории поля, где вытекающие из наличия группы симметрии законы сохранения часто оказываются единственным источником информации о свойствах системы. Для формального вывода Н. т. в (классической или квантовой) теории поля рассматривают интеграл действия:



\[

S=\int_{R} L\left[\varphi^{a}(x), \varphi_{, v}^{a}(x) ; x^{\mu}\right] d^{4} x ; \quad \mu, v=0,1,2,3,

\]



где $L\left[\varphi^{a}(x), \varphi_{, v}^{a}(x) ; x^{\mu}\right]$ - лагранжиан, зависящий от функций поля $\varphi^{a}(x)$, их первых производных по всем четырем координатам $\varphi_{, v}^{a} \equiv \partial \varphi^{a} / \partial x^{\nu}$ и, возможно, от координат $x^{\mu}$ ( $x=\left\{x^{\mu}\right\}-$ точка пространства-времени; индекс $a$ нумерует





\end{document}